% LNCS format - Springer Lecture Notes in Computer Science
\documentclass[runningheads]{llncs}

\usepackage[T1]{fontenc}
\usepackage{lmodern}
\usepackage{graphicx}
\usepackage{booktabs}
\usepackage{amsmath}
\usepackage{multirow}
\usepackage{hyperref}
\usepackage{xcolor}
\usepackage{subcaption}
\usepackage{float}
\usepackage{placeins}

\pagestyle{plain}
\raggedbottom
\widowpenalty=10000
\clubpenalty=10000

\begin{document}

\title{Fair-Share Scheduling for Multi-Tenant\\Serverless Platforms}

\author{COMP 6910 --- Joshua Oseimobor (202492785) , Isaac Antwi Okraku (202580677)}
\authorrunning{J. Oseimobor, I. Antwi Okraku}

\institute{Memorial University of Newfoundland Labrador}

\maketitle

\begin{abstract}
Serverless platforms promise all tenants the same baseline service-level agreement (SLA), yet current schedulers allow large tenants to starve smaller ones through disproportionate resource consumption. We present a two-phase fairness-guarantee scheduler that structurally prevents starvation by guaranteeing each tenant a minimum dispatch share before filling remaining capacity efficiently. Using a discrete-event simulator built with SimPy, we evaluate our scheduler against three baselines---FIFO, Round-Robin, and Shortest-Job-First (SJF)---across five experiment scenarios with 25 tenants, 8 servers, and approximately 39,000 invocations per run. Our scheduler achieves the highest resource fairness (Jain's Index 0.305--0.309 vs.\ SJF's 0.216--0.300) while ensuring zero SLA violations for small and medium tenants across all experiments. We show that SJF's apparent fairness is accidental---it benefits small tenants only when their workloads happen to be lightweight---whereas our scheduler provides principled, tenant-aware protection regardless of workload characteristics.

\keywords{Serverless computing \and Multi-tenant fairness \and Resource scheduling \and SLA enforcement \and Discrete-event simulation}
\end{abstract}

%----------------------------------------------------------------------
\section{Introduction}\label{sec:intro}
%----------------------------------------------------------------------

Serverless computing platforms such as AWS Lambda, Azure Functions, and Google Cloud Functions abstract away infrastructure management, promising tenants consistent performance under a shared SLA. Large-scale workload studies confirm that serverless invocation frequencies span eight orders of magnitude across tenants, with most functions executing in under one second~\cite{shahrad2020serverless}. However, this abstraction masks a fundamental fairness problem: when multiple tenants share finite compute resources, current schedulers do not enforce proportional allocation. A large tenant submitting heavy or bursty workloads can consume disproportionate resources, pushing smaller tenants into longer queues, more cold starts, and degraded response times---effectively violating their SLA even though the platform remains nominally available.

Cold starts compound this problem. Production measurements show that cold start latency ranges from tens of milliseconds to several seconds depending on runtime and platform architecture~\cite{wang2018peeking}. Tenants deprioritized by the scheduler are more likely to land on servers without warm containers, incurring additional latency penalties. Recent work on reducing cold starts addresses the problem in isolation but does not consider the fairness implications of which tenants bear this burden.

We address the following research question: \emph{Can a lightweight, adaptive scheduler protect vulnerable tenants' baseline SLA guarantees in a shared serverless environment---preventing large tenants from starving smaller ones---while maintaining fair resource distribution and acceptable overall throughput under bursty, heterogeneous workloads?}

Our contributions are:
\begin{enumerate}
    \item We identify and demonstrate the starvation problem with baseline schedulers (FIFO, Round-Robin, SJF) in multi-tenant serverless environments through controlled simulation experiments.
    \item We propose a two-phase scheduler that guarantees each tenant a minimum dispatch share per scheduling round (Phase~1: fairness) before filling remaining capacity via global SJF (Phase~2: efficiency), providing principled tenant-aware protection with low-millisecond overhead.
    \item We provide a reproducible simulation framework with comprehensive evaluation across five experiment types, showing that our scheduler protects small and medium tenants' SLA while maintaining the highest resource fairness.
\end{enumerate}

%----------------------------------------------------------------------
\section{Related Work}\label{sec:related}
%----------------------------------------------------------------------

\paragraph{Baseline Schedulers.}
FIFO processes invocations in arrival order, causing head-of-line blocking when large tenants flood the queue. Round-Robin cycles through tenants equally, but equal turns do not equate to equal resources when function costs vary---Mace et al.~\cite{mace20162dfq} show that one-dimensional fair queuing is insufficient when request costs differ. SJF minimizes average latency by prioritizing short functions but systematically starves tenants with heavier workloads and has no awareness of tenants, SLAs, or fairness. Dominant Resource Fairness (DRF)~\cite{ghodsi2011dominant} allocates based on each tenant's dominant resource demand with desirable theoretical properties, but was designed for long-running batch systems and cannot handle sub-second serverless invocations, cold start penalties, or event-driven traffic bursts.

\paragraph{Serverless Scheduling.}
Recent work demonstrates that scheduling choices significantly impact serverless performance and cost, with the default Linux CFS scheduler shown to be suboptimal for short-lived, bursty invocations. However, existing serverless scheduling work optimizes single-tenant performance, not multi-tenant resource sharing.

\paragraph{Multi-Tenant Fairness.}
2DFQ~\cite{mace20162dfq} introduces two-dimensional fair queuing that accounts for both request rate and cost, demonstrating that effective fairness requires resource-aware, multi-dimensional mechanisms. However, existing multi-tenant fairness systems target storage and general cloud services, not the unique sub-second timescales and cold-start dynamics of serverless.

\paragraph{Novelty Gap.}
No existing scheduler addresses all four problems simultaneously in a serverless context: (1) head-of-line blocking from large tenants, (2) resource heterogeneity across function types, (3) starvation of tenants with heavier workloads, and (4) the inability of batch-oriented fairness schemes to operate at serverless timescales with cold starts and bursts.

%----------------------------------------------------------------------
\section{System Model}\label{sec:model}
%----------------------------------------------------------------------

\paragraph{Tenant Model.}
We model 25 tenants as distinct organizations sharing a serverless platform, categorized into three tiers: 5 large tenants (50--100 invocations/sec, 60\% heavy functions), 10 medium tenants (10--30 inv/sec, mixed workloads), and 10 small tenants (1--5 inv/sec, 70\% lightweight functions). All tenants share the same baseline SLA.

\paragraph{Infrastructure.}
The platform comprises 8 homogeneous servers, each with 4,000 millicores CPU and 8,192~MB memory, totaling 32,000m CPU and 65,536~MB cluster-wide. Function invocations are classified into three archetypes: \emph{lightweight} (100m CPU, 64~MB, 5--20ms), \emph{medium} (250m CPU, 256~MB, 30--80ms), and \emph{heavy} (500m CPU, 512~MB, 100--300ms).

\paragraph{Cold Start Model.}
Informed by production measurements~\cite{wang2018peeking}, each server maintains warm containers with a 300-second TTL. Cold starts add a fixed +50ms penalty. All containers start cold at simulation initialization.

\paragraph{SLA Guarantees.}
All tenants share a uniform SLA: P95 response time below 1,000ms and throughput at or above 80\% of the submitted invocation rate. The fairness problem is that baseline schedulers allow large tenants to push smaller tenants below these thresholds.

\paragraph{Fairness Definition.}
We adopt equal proportional fairness: each active tenant is entitled to an equal share of cluster resources ($\text{total\_capacity} / n_{\text{active}}$). No tenant should be starved below their SLA guarantees because another tenant consumes disproportionate resources. Fairness is measured using Jain's Fairness Index~\cite{jain1984quantitative} over the ratio of actual resource consumption to equal entitlement.

%----------------------------------------------------------------------
\section{Proposed Approach}\label{sec:approach}
%----------------------------------------------------------------------

\subsection{Two-Phase Fair-Share Scheduler}

Rather than computing a single priority score, our scheduler operates in two distinct phases per scheduling round, structurally guaranteeing fairness before optimizing efficiency.

\paragraph{Phase 1: Fairness Guarantee.}
Each active tenant with pending invocations receives exactly one dispatch per scheduling round. Tenants are ordered by their \emph{dispatch deficit}---the difference between the tenant's proportional share and its actual dispatch count within the current sliding window:
\begin{equation}\label{eq:deficit}
    \text{Deficit}_i = \frac{\sum_{j} D_j}{n_{\text{active}}} - D_i
\end{equation}
where $D_i$ is the number of invocations dispatched for tenant~$i$ in the current window, $\sum_j D_j$ is the total dispatches across all tenants, and $n_{\text{active}}$ is the number of tenants with pending work. Tenants with the highest deficit (most under-served) are scheduled first. Within each tenant's queue, the shortest job is selected (intra-tenant SJF), combining fairness across tenants with efficiency within each tenant.

\paragraph{Phase 2: Efficiency.}
After every active tenant has received its guaranteed dispatch, remaining server capacity is filled using global SJF across all pending invocations regardless of tenant. This maximizes throughput utilization without compromising the fairness floor established in Phase~1.

\paragraph{Sliding Window.}
Dispatch counts are tracked in a configurable sliding window (default 5.0 seconds). When the window expires, all counts reset, preventing historical imbalances from permanently biasing scheduling decisions. The ablation study (Section~\ref{sec:ablation}) validates this window size.

\paragraph{Server Selection.}
Invocations are assigned to servers using a two-step preference: (1)~servers with a warm container for the function type and sufficient capacity, avoiding cold start penalties; (2)~among candidates, the least-loaded server by CPU utilization, with memory as tiebreaker.

\paragraph{Design Rationale.}
The two-phase design prevents starvation \emph{structurally} rather than \emph{reactively}. Unlike weighted-score approaches that must tune parameters to balance fairness and efficiency, our scheduler guarantees a fairness floor (Phase~1) and then maximizes efficiency with remaining capacity (Phase~2). Over-consuming tenants are naturally constrained: they receive their one guaranteed dispatch per round, but Phase~2 capacity is allocated by job length, not tenant identity---so heavy functions from large tenants wait behind lighter work.

%----------------------------------------------------------------------
\section{Implementation}\label{sec:implementation}
%----------------------------------------------------------------------

We implement a discrete-event simulator using Python and SimPy~\cite{mueller2024simpy}, following the established methodology of evaluating cloud scheduling algorithms via simulation~\cite{calheiros2011cloudsim}. Synthetic workloads are parameterized based on real serverless workload characteristics from the Azure Functions Trace~\cite{shahrad2020serverless}.

\paragraph{Simulation Engine.}
Three concurrent SimPy processes drive the simulation: (1)~an \emph{arrival process} that feeds invocations into tenant pending queues at their arrival times; (2)~a \emph{scheduler tick process} that runs every 10ms simulation time, invoking the two-phase scheduler to batch arrivals and completions between ticks; and (3)~an \emph{execution process} (one per assigned invocation) that reserves resources, simulates execution, and releases resources on completion. A fourth background process evicts expired warm containers every second. The tick-based architecture reduces scheduling calls from approximately 72,000 (per-event) to approximately 6,000 per 60-second run, cutting experiment runtime from 30~minutes to 3~minutes while more closely modeling real scheduler behavior.

\paragraph{Workload Generation.}
For each tenant, invocations arrive as a Poisson process at the tenant's configured rate. Function type is sampled from the tenant's profile distribution, resource demands are deterministic per archetype, and durations are uniformly sampled within archetype ranges. Burst experiments inject a configurable multiplier to a specified tenant's arrival rate during a time window.

\paragraph{Reproducibility.}
All experiments use fixed random seeds, version-controlled YAML configurations, and deterministic SimPy scheduling. The full codebase, raw data, processed metrics, and configuration files are included in the project submission.

%----------------------------------------------------------------------
\section{Evaluation}\label{sec:evaluation}
%----------------------------------------------------------------------

We evaluate our Fair-Share scheduler against three baselines---FIFO, Round-Robin, and SJF---across five experiment scenarios. Each 60-second run generates approximately 39,000 invocations across 25 tenants and 8 servers.

\subsection{Metrics}

We use Jain's Fairness Index~\cite{jain1984quantitative} as our primary fairness metric. Given a vector of per-tenant values $\mathbf{x} = (x_1, \ldots, x_n)$, Jain's Index is:
\begin{equation}\label{eq:jains}
    J(\mathbf{x}) = \frac{\big(\sum_{i=1}^{n} x_i\big)^2}{n \cdot \sum_{i=1}^{n} x_i^2}
\end{equation}
where $J = 1.0$ indicates perfect equality and lower values indicate increasing disparity. We apply this formula to two distinct input vectors to capture complementary dimensions of fairness:

\begin{itemize}
    \item \textbf{Jain's Resource Fairness} (primary): each $x_i$ is the tenant's \emph{fair share ratio}---the ratio of actual CPU-seconds consumed to proportional entitlement ($\text{total CPU-seconds} / n_{\text{active}}$). This measures whether all tenants receive equitable resource allocation.
    \item \textbf{Jain's SLA Compliance} (secondary): each $x_i = 1 - v_i$, where $v_i$ is the tenant's invocation-level SLA violation rate. This measures whether all tenants receive equitable service quality---a scheduler could have low overall violations but concentrate them on a few tenants.
    \item \textbf{SLA Violation Rate}: reported at three levels---per-tenant binary (P95 $>$ 1,000ms or throughput $<$ 80\%), tenant fraction, and invocation-weighted.
    \item \textbf{P95 Latency}: per tenant size category.
    \item \textbf{Scheduling Overhead}: wall-clock time per scheduling decision.
\end{itemize}

\subsection{Experiment 1: Steady-State Fairness}

All tenants operate at constant rates for 60 seconds. Table~\ref{tab:steady} summarizes the results.

\begin{table}[!ht]
\centering
\caption{Steady-state experiment results. ``Tenants Violating'' shows the count of small/medium/large tenants exceeding SLA thresholds.}\label{tab:steady}
\small
\begin{tabular}{@{}lcccc@{}}
\toprule
\textbf{Metric} & \textbf{FIFO} & \textbf{RR} & \textbf{SJF} & \textbf{Fair-Share} \\
\midrule
Jain's Resource Fairness & 0.299 & 0.306 & 0.300 & \textbf{0.305} \\
Jain's SLA Compliance & 0.944 & 0.885 & \textbf{1.000} & 0.941 \\
Invocation Violation Rate & 0.899 & 0.388 & \textbf{0.029} & 0.344 \\
Tenant Violation Rate & 1.00 & 0.12 & \textbf{0.00} & 0.20 \\
\midrule
Small Tenants Violating & 10/10 & 0/10 & 0/10 & 0/10 \\
Medium Tenants Violating & 10/10 & 0/10 & 0/10 & 0/10 \\
Large Tenants Violating & 5/5 & 3/5 & 0/5 & 5/5 \\
\midrule
Overhead (ms, avg) & 0.46 & 0.06 & 0.22 & 0.74 \\
\bottomrule
\end{tabular}
\end{table}

FIFO collapses under load---all 25 tenants violate their SLA because large tenants' invocations block the queue. Round-Robin and Fair-Share both protect small and medium tenants (0 violations), but Fair-Share achieves marginally higher resource fairness. SJF achieves the lowest overall violation rate and highest SLA compliance, but this warrants careful interpretation (see Section~\ref{sec:discussion}).

\subsection{Experiment 2: Burst Resilience}

A large tenant spikes to 5$\times$ its normal arrival rate for 10 seconds at $t=20$s. Table~\ref{tab:burst} shows the results.

\begin{table}[!ht]
\centering
\caption{Burst test results. A large tenant spikes 5$\times$ for 10 seconds.}\label{tab:burst}
\small
\begin{tabular}{@{}lcccc@{}}
\toprule
\textbf{Metric} & \textbf{FIFO} & \textbf{RR} & \textbf{SJF} & \textbf{Fair-Share} \\
\midrule
Jain's Resource Fairness & 0.254 & \textbf{0.307} & 0.268 & 0.306 \\
Invocation Violation Rate & 0.710 & 0.359 & \textbf{0.056} & 0.337 \\
\midrule
Small Violating & 10/10 & 0/10 & 0/10 & 0/10 \\
Medium Violating & 10/10 & 0/10 & 0/10 & 0/10 \\
Large Violating & 5/5 & 3/5 & 5/5 & 5/5 \\
\bottomrule
\end{tabular}
\end{table}

FIFO again collapses entirely. Under burst conditions, even SJF's large tenants all violate SLA (5/5), unlike steady-state where none did. Our scheduler absorbs the burst without pushing small or medium tenants into violation, demonstrating that its fairness protection is robust to sudden load changes.

\subsection{Experiment 3: Skewed Load}

Two large tenants are configured to submit 200--250 invocations/sec each, generating approximately 70\% of all traffic. This is the most adversarial scenario for fairness.

\begin{table}[!ht]
\centering
\caption{Skewed load results---the most adversarial fairness scenario.}\label{tab:skewed}
\small
\begin{tabular}{@{}lcccc@{}}
\toprule
\textbf{Metric} & \textbf{FIFO} & \textbf{RR} & \textbf{SJF} & \textbf{Fair-Share} \\
\midrule
Jain's Resource Fairness & 0.207 & \textbf{0.309} & 0.216 & \textbf{0.309} \\
Invocation Violation Rate & 0.990 & 0.498 & \textbf{0.021} & 0.446 \\
\midrule
Small Violating & 10/10 & 0/10 & 0/10 & 0/10 \\
Medium Violating & 10/10 & 0/10 & 0/10 & 0/10 \\
Large Violating & 5/5 & 4/5 & 1/5 & 5/5 \\
\bottomrule
\end{tabular}
\end{table}

Under extreme skew, the resource fairness gap widens significantly: Fair-Share achieves Jain's Index 0.309 versus SJF's 0.216---a 43\% improvement. SJF provides the \emph{most unfair resource distribution} of any scheduler despite low overall violation rates. FIFO nearly completely collapses (99\% invocation violation rate). Both Round-Robin and Fair-Share protect all small and medium tenants while achieving the highest resource fairness.

\subsection{Experiment 4: Scalability}

We run the Fair-Share scheduler at default configuration (25 tenants, 8 servers) and measure scheduling overhead. The average scheduling overhead is 0.74--8.06ms depending on load, with P95 at 1.09--15.34ms. The overhead scales linearly with tenant count and remains in the low-millisecond range, confirming the scheduler is practical for real-time operation.

\subsection{Experiment 5: Ablation Study}\label{sec:ablation}

We sweep the sliding window duration to validate the default value of 5.0 seconds. The sliding window controls how frequently dispatch counts reset---shorter windows react faster to load changes but may over-correct, while longer windows provide more stable fairness tracking but respond slowly to bursts. Key observations from the ablation data:
\begin{itemize}
    \item Very short windows (e.g., 1.0s) cause frequent resets that prevent the deficit from accumulating meaningful signal, reducing the scheduler's ability to detect and correct sustained imbalances.
    \item Very long windows (e.g., 30.0s) allow historical dispatch imbalances to persist too long, delaying correction when workload patterns shift mid-experiment.
    \item The default 5.0-second window balances responsiveness with stability: deficit tracking is long enough to capture sustained starvation patterns while resetting frequently enough to adapt to changing conditions.
\end{itemize}

%----------------------------------------------------------------------
\section{Discussion}\label{sec:discussion}
%----------------------------------------------------------------------

\paragraph{Why Not SJF?}
A natural question is why not simply use SJF, given its low overall SLA violation rates (2.1--5.6\% across experiments). The answer is that SJF's low latency for small tenants is \emph{accidental}, not principled. SJF prioritizes short functions---and small tenants happen to submit mostly lightweight functions. But SJF has no awareness of tenants, SLAs, or fairness. Under skewed load, SJF produces the most unfair resource distribution of any scheduler (Jain's 0.216 vs.\ Fair-Share's 0.309). In a real platform where any tenant might have heavy functions, SJF would silently starve them with no mechanism to detect or correct it. Our scheduler provides \emph{principled} protection: it guarantees each tenant a minimum dispatch share per scheduling round, then fills remaining capacity efficiently---regardless of function characteristics.

\paragraph{Large Tenant Penalty as Design Feature.}
In all experiments, our scheduler's large tenants show high SLA violation rates. This is a deliberate consequence of finite resources: when resources are constrained, protecting smaller tenants' baseline SLA necessarily means that over-consuming tenants bear proportional cost. This mirrors real-world platform behavior---AWS Lambda enforces per-account concurrency limits for the same reason. The key property is that the penalty is \emph{proportional} (larger consumers absorb more degradation) rather than \emph{arbitrary} (as with FIFO, where all tenants suffer regardless of consumption).

\paragraph{Trade-Off Analysis.}
Table~\ref{tab:tradeoff} summarizes the fundamental trade-off across schedulers. No scheduler dominates all dimensions simultaneously, which is expected: fairness and raw throughput are inherently in tension under finite resources. Fair-Share and Round-Robin occupy the same region of the Pareto frontier for resource fairness and tenant protection, but Fair-Share adds deficit-aware ordering that Round-Robin lacks---prioritizing the most under-served tenants rather than cycling blindly.

\begin{table}[!ht]
\centering
\caption{Trade-off summary across all experiments.}\label{tab:tradeoff}
\small
\begin{tabular}{@{}lccc@{}}
\toprule
\textbf{Scheduler} & \textbf{Resource Fairness} & \textbf{Small/Med Protection} & \textbf{Overall Violation} \\
\midrule
FIFO & Low (0.207--0.299) & None & Very High \\
Round-Robin & High (0.306--0.309) & Full & High \\
SJF & Low (0.216--0.300) & Accidental & Very Low \\
Fair-Share & High (0.305--0.309) & Principled & Moderate \\
\bottomrule
\end{tabular}
\end{table}

%----------------------------------------------------------------------
\section{Limitations and Future Work}\label{sec:limitations}
%----------------------------------------------------------------------

\paragraph{Simulation vs.\ Real Deployment.}
Our evaluation uses a discrete-event simulator rather than a real cloud platform. While SimPy accurately models scheduling logic and resource contention, it cannot capture real-world effects such as network latency, OS scheduling interference, JVM/runtime warm-up, or hardware heterogeneity.

\paragraph{Fixed Cold Start Penalty.}
We model cold starts as a fixed +50ms penalty. In production, cold start latency varies significantly with runtime, container size, and platform architecture~\cite{wang2018peeking}---from tens of milliseconds for lightweight runtimes to several seconds for heavy ones.

\paragraph{No Network Latency.}
The simulator assumes negligible network latency between the scheduler and servers. In geo-distributed deployments, network delays could affect scheduling decisions and their effectiveness.

\paragraph{Future Directions.}
(1)~Deployment on a real cloud platform (e.g., Kubernetes with Knative) to validate simulation findings. (2)~Tiered SLA extensions (Gold/ Silver /Bronze) where tenants pay for different service guarantees. (3)~Cost modeling to quantify the financial impact of fairness-aware scheduling. (4)~Variable cold start penalties calibrated to production measurements.

%----------------------------------------------------------------------
\section{Conclusion}\label{sec:conclusion}
%----------------------------------------------------------------------

Shared serverless platforms allow large tenants to starve smaller ones under contention---a fairness problem that current schedulers do not address. We proposed a two-phase fairness-guarantee scheduler that structurally prevents starvation by guaranteeing each tenant a minimum dispatch share before filling remaining capacity efficiently. Through comprehensive simulation across five experiment types, we demonstrated that our scheduler achieves the highest resource fairness (Jain's Index 0.305--0.309) while ensuring zero SLA violations for small and medium tenants. We showed that SJF's apparently superior overall metrics mask fundamentally unfair resource distribution---its protection of small tenants is accidental, not principled. Our two-phase approach provides a practical, principled alternative: low-millisecond scheduling overhead, structural fairness guarantees, and transparent trade-offs where over-consuming tenants bear proportional cost.

All source code, experiment configurations, raw data, and processed metrics are included in the project repository for full reproducibility.

%----------------------------------------------------------------------
% References
%----------------------------------------------------------------------
\begin{thebibliography}{7}

\bibitem{wang2018peeking}
Wang, L., Li, M., Zhang, Y., Ristenpart, T., Swift, M.: Peeking behind the curtains of serverless platforms. In: USENIX ATC '18 (2018)

\bibitem{shahrad2020serverless}
Shahrad, M., Fonseca, R., Goiri, I., et al.: Serverless in the wild: Characterizing and optimizing the serverless workload at a large cloud provider. In: USENIX ATC '20 (2020)

\bibitem{ghodsi2011dominant}
Ghodsi, A., Zaharia, M., Hindman, B., Konwinski, A., Shenker, S., Stoica, I.: Dominant resource fairness: Fair allocation of multiple resource types. In: NSDI '11 (2011)

\bibitem{mace20162dfq}
Mace, J., Bodik, P., Fonseca, R., Musuvathi, M.: 2DFQ: Two-dimensional fair queuing for multi-tenant cloud services. In: SoCC '16 (2016)

\bibitem{jain1984quantitative}
Jain, R.K., Chiu, D.-M.W., Hawe, W.R.: A quantitative measure of fairness and discrimination for resource allocation in shared computer systems. DEC Research Report TR-301 (1984)

\bibitem{calheiros2011cloudsim}
Calheiros, R.N., Ranjan, R., Beloglazov, A., De Rose, C.A.F., Buyya, R.: CloudSim: A toolkit for modeling and simulation of cloud computing environments. Software: Practice and Experience \textbf{41} (2011)

\bibitem{mueller2024simpy}
Mueller, K.: Discrete event simulation: It's easy with SimPy! arXiv:2405.01562 (2024)

\end{thebibliography}

\end{document}
